\section{ Research Experience }
\subsection{Boston University}
    \cventry{Jan 2023 -- Present}%
      {\textbf{Advisor(s):} {\href{https://www.batman-lab.com/} { \textbf{\textcolor{blue}{Dr. Kayhan Batmanghelich}}}},%
      {\href{https://www.bumc.bu.edu/camed/profile/clare-poynton/}
        { \textbf{\textcolor{blue}{Dr. Clare B. Poynton}}}}}
      {Graduate Research Assistant}%
      {Boston, Massachusetts, USA}%
      {}%
      {
        \begin{itemize}[itemsep=0pt, parsep=0pt]
        % \item {\bf Advisor(s):} {\href{https://www.batman-lab.com/}
        % { \textbf{\textcolor{blue}{Dr. Kayhan Batmanghelich}}}}
        % \item {\bf Research Area}: Explainable AI; Computer Vision; Medical Imaging.
        % \item Currently extending Mammo-CLIP to develop the $1^{st}$ generative foundation model using screening mammograms for breast, capable of diagnosing breast cancer (diagnosis), predicting risk of having breast cancer (prognosis), and automated report generation.
        \item Developed \textbf{LADDER}, a slice discovery and mitigation algorithm using vision language (VLM) models and LLMs to reason and fix the classifier's mistakes. {\textbf{Accepted at ACL, 2025 (Findings)}}
        \item Developed \textbf{Mammo-CLIP} the first vision language foundation model for 2D mammograms.  \href{https://shantanu-ai.github.io/projects/MICCAI-2024-Mammo-CLIP/}{\textbf{Accepted at MICCAI, 2024 (top 11\%)}}
        \item Applied the mixture of interpretable models for efficient transfer learning to an unseen domain with limited training data.
        \href{https://shantanu-ai.github.io/projects/MICCAI-2023-MoIE-CXR/}{\textbf{Accepted at MICCAI, 2023 (top 14\%)}}.
        \end{itemize}
      }
  
    \subsection{University of Pittsburgh}
    \cventry{Aug 2021 -- Dec 2022}%
      {\textbf{Advisor(s):} {\href{https://www.batman-lab.com/} { \textbf{\textcolor{blue}{Dr. Kayhan Batmanghelich}}}},
      {\href{https://forougha.github.io/}
        { \textbf{\textcolor{blue}{Dr. Forough Arabshahi}}}}}%
      {Graduate Student Researcher}%
      {Pittsburgh, Pennsylvania, USA}%
      {}%
      {
        \begin{itemize}[itemsep=0pt, parsep=0pt]
        % \item {\bf Advisor(s):} 
        % {\href{https://www.batman-lab.com/}
        % { \textbf{\textcolor{blue}{Dr. Kayhan Batmanghelich}}}}, {\href{https://forougha.github.io/}
        % { \textbf{\textcolor{blue}{Dr. Forough Arabshahi}}}} 
        % \item {\bf Research Area}: Explainable AI; Computer Vision; Medical Imaging.
        \item Developed an iterative algorithm to extract a mixture of interpretable models from a Blackbox to provide instance-specific First-order logic-based explanations using human-understandable concepts. \href{https://shantanu-ai.github.io/projects/ICML-2023-MoIE/}
        {\textbf{Accepted at ICML, 2023}}.
        % \item Localized Pneumonia and Pneumothorax from \textbf{MIMIC-CXR} dataset by leveraging the anatomical landmarks (weak labels) using the \textbf{Stanford RadGraph NLP pipleline}. 
        % Also, designed the baseline using \textbf{RetinaNet}.
        % \href{https://conferences.miccai.org/2022/papers/045-Paper1726.html}
        % {\textbf{Accepted at MICCAI, 2022 and RAD: AI, 2024}}.
        % \item Investigated why \textbf{lottery ticket hypothesis} works using: \textbf{Concept activation vectors (TCAV)} and  \textbf{Grad-CAM}. 
        % \href{https://github.com/shantanu-ai/Explainability-with-LTH}
        % {\textcolor{blue}{\textbf{[Code]}}} \href{https://github.com/shantanu-ai/Explainability-with-LTH/blob/main/doc/VLR.pdf}
        % {\textcolor{blue}{\textbf{[Report]}}}
        \end{itemize} 
      }

      % 1. Currently developing a novel algorithm to explain the prediction of a black box classifier with the help of network pruning using lottery ticket hypothesis. In this project, we aim to iteratively extract the explainable concepts from a black-box model using a neuro-symbolic explainable model. Also, we aim to detect the shortcut (biased) concepts from the black box and iteratively remove them from the features of the black box to make it robust.
      % 2. Analyzed why pruning strategies using lottery ticket hypothesis works or fails in terms of explainability metrics -- Concept activation vectors (TCAV) for global concepts and saliency maps like Grad-CAM, Integrated-Gradient for pixel attributions.
      % 3. Developed an attention model to leverage the anatomical landmarks (weak labels) from Stanford RadGraph NLP pipleline to detect Pneumonia and Pneumothorax from MIMIC-CXR dataset. Also, designed the baseline using RetinaNet. 
  
    \subsection{University of Florida}
    \cventry{Jan 2020 -- Jul 2021}%
      {{\bf Advisor(s):} 
        {\href{https://www.cise.ufl.edu/~butler/}
        { \textbf{\textcolor{blue}{Dr. Kevin Butler}}}},
        {\href{https://scholar.google.com/citations?user=ysr--voAAAAJ&hl=en}
        { \textbf{\textcolor{blue}{Dr. Jiang Bian}}}},
        {\href{https://epidemiology.phhp.ufl.edu/profile/prosperi-mattia/}
        { \textbf{\textcolor{blue}{Dr. Mattia Prosperi}}}}
        }%
      {Graduate Research Assistant}%
      {Gainesville, Florida, USA}%
      {}%
      {
       \begin{itemize}[itemsep=0pt, parsep=0pt]
        % \item {\bf Research Area}: Causal Inference, Deep Learning.
        \item Developed novel deep learning frameworks to estimate propensity score, namely DPN-SA (\textbf{JAMIA 2021}), PSSAM-GAN (\textbf{CMPB-U 2021}), and DR-VIDAL (\textbf{AMIA 2022, oral}), for the efficient estimation of individual treatment effects (ITE).
        \end{itemize}
      }
        
      % {
      %   \begin{itemize}[itemsep=0pt, parsep=0pt]
      %   % \item {\bf Advisor(s):} {\href{https://epidemiology.phhp.ufl.edu/profile/prosperi-mattia/}
      %   % { \textbf{\textcolor{blue}{Dr. Mattia Prosperi}}}}, 
      %   % {\href{https://www.cise.ufl.edu/~butler/}
      %   % { \textbf{\textcolor{blue}{Dr. Kevin Butler}}},
      %   % }
      %   % \item {\bf Research Area}: Causal Inference, Deep Learning.
      %   \item Developed a novel deep learning framework to (1) generate the counterfactual outcomes based on treatment using a Generative Adversarial Network with \textbf{information-theoretic} regularization; (2) utilized the counterfactual outcomes to estimate the individual treatment effect (\textbf{ITE}) using \textbf{doubly robust optimization} for faster convergence. 
      %   \href{https://www.ncbi.nlm.nih.gov/pmc/articles/PMC10148269/}
      %   {\textbf{Accepted at AMIA Symposium (Oral), 2022}}.
      %   \item Designed a novel algorithm using a Generative Adversarial Network to generate synthetic treated samples to remove imbalance within an observational dataset for \textbf{Propensity score matching}. \href{https://www.sciencedirect.com/science/article/pii/S2666990021000197/}
      %   {\textbf{Accepted at Computer Methods and Programs in Bio-medicine Update}}.
      %   \item Developed a \textbf{sparse autoencoder} to reduce the dimensionality of the covariates of the patients to calculate the \textbf{Propensity score} in an efficient way to estimate the average treatment effect (\textbf{ATE}) of the treatment. \href{https://academic.oup.com/jamia/article-abstract/28/6/1197/6139936}
      %   {\textbf{Accepted at JAMIA}.}
      %   \end{itemize}
      % }

      %    1. Designed a robust deep learning model amalgamating the theories of Causal Graphs and Deep Variational Information Bottleneck. Studied the failure modes of the causal model by performing adversarial attacks.
      %    2. Developed the baseline using the experiments in the papers "Deep Variational Information Bottleneck" (Alemi et al.) and "A Causal View on Robustness of Neural Networks" (Zhang et al.) in Pytorch.
  
    % \subsection{University of Florida}
    % \cventry{Jan 2020 -- Feb 2021}%
    %   {Data Intelligence Systems Lab (DISL)}%
    %   {Research Assistant}%
    %   {Gainesville, Florida, USA}%
    %   {}%
    %   {
    %     \begin{itemize}[itemsep=0pt, parsep=0pt]
    %     \item {\bf Advisor(s):} {\href{https://epidemiology.phhp.ufl.edu/profile/prosperi-mattia/}
    %     { \textbf{\textcolor{blue}{Dr. Mattia Prosperi}}}},
    %     {\href{https://scholar.google.com/citations?user=ysr--voAAAAJ&hl=en}
    %     { \textbf{\textcolor{blue}{Dr. Jiang Bian}}}} 
    %     % \item {\bf Research Area}: Causal Inference, Deep Learning.
    %     \item Designed a novel algorithm using a Generative Adversarial Network to generate synthetic treated samples to remove imbalance within an observational dataset for \textbf{Propensity score matching} [J2]. \href{https://www.sciencedirect.com/science/article/pii/S2666990021000197/}
    %     {\textbf{Accepted at Computer Methods and Programs in Bio-medicine Update}}.
    %     \item Developed a \textbf{sparse autoencoder} to reduce the dimensionality of the covariates of the patients to calculate the \textbf{Propensity score} in an efficient way to estimate the average treatment effect (\textbf{ATE}) of the treatment [J1]. \href{https://academic.oup.com/jamia/article-abstract/28/6/1197/6139936}
    %     {\textbf{Accepted at JAMIA}}
    %     \end{itemize}
    %   }